% !TeX program = xelatex
% In Overleaf, select XeLaTeX in the project Compiler settings.
\documentclass[UTF8,fontset=fandol,11pt,a4paper]{ctexart}
\usepackage[margin=24mm]{geometry}
\usepackage{amsmath,amssymb,booktabs,enumitem}
\usepackage[colorlinks=true,linkcolor=blue,urlcolor=blue]{hyperref}
\setlist[enumerate]{itemsep=3pt,topsep=5pt}
\setlength{\emergencystretch}{2em}
\setcounter{secnumdepth}{0}
\newcommand{\bits}{\{0,1\}}
\newcommand{\NAND}{\operatorname{NAND}}
\newcommand{\MAJ}{\operatorname{MAJ}}
\newcommand{\LOOKUP}{\operatorname{LOOKUP}}
\newcommand{\SIZE}{\operatorname{SIZE}}
\title{理论计算机科学导引\\第 2、3、5 章选题解答（中英双语）\\
\large Selected Exercise Solutions: Chapters 2, 3, and 5}
\author{}
\date{}

\begin{document}
\maketitle
\vspace{-2em}
\noindent
依据所提供的 \texttt{lnotes\_book.pdf}，解答习题
2.4、2.5(1)、2.14、3.9、3.14、5.1、5.2、5.3、5.4、5.9。
题目所在页分别为 116、118、156--157、211--212 页。
以下所有 $\log$ 均以 $2$ 为底，$[n]=\{0,\ldots,n-1\}$；
电路大小指 NAND 门数，允许扇出，连线本身不计门数。

\medskip
\noindent\textbf{特别说明：}
5.4 原文对 $m$ 没有限制，这一表述并不成立。
该题先给出反例，再证明 $1\le m\le 2^n$ 范围内的正确版本。
这是对原题量词的检查，不是教材中已经声明的附加条件。

\medskip
\noindent
This document contains complete Chinese and English solutions to Exercises
2.4, 2.5(1), 2.14, 3.9, 3.14, 5.1, 5.2, 5.3, 5.4, and 5.9.
The English solutions form the second part. Exercise 5.4 is false as stated
when the output length is unrestricted; both parts include a counterexample
and a proof of a corrected statement.

\tableofcontents

\clearpage
\phantomsection
\addcontentsline{toc}{part}{中文解答}
\part*{中文解答}

\section{习题 2.4：有界度图的短编码}

\paragraph{题意}
令 $G_n$ 为顶点集为 $[n]$、无自环、每个顶点度数至多 $10$ 的有向图集合。
证明：充分大的 $n$ 存在单射
\[
 E:G_n\longrightarrow\bits^{\lfloor1000n\log n\rfloor}.
\]

\paragraph{构造}
令 $\ell=\lceil\log n\rceil$。按顶点 $v=0,1,\ldots,n-1$ 的顺序，
记录 $v$ 的出邻接表。因为出度不超过 $10$，每一行都可写为：
\begin{enumerate}
 \item 用 $4$ 位记录出度 $d_v\in\{0,1,\ldots,10\}$；
 \item 将出邻居按编号递增排列，每个编号用 $\ell$ 位表示，共预留 $10$ 个位置；
 \item 不足 $10$ 个的位置补全零。
\end{enumerate}
注意：编号 $0$ 本身是合法邻居，但不会和补零混淆，因为 $d_v$ 指定了有效位置数。
无论题目中的度数指总度数还是出度，上述出度上界都成立。

原始编码长度为
\[
 L=n\bigl(4+10\lceil\log n\rceil\bigr).
\]
当 $n\ge2$ 时，$\lceil\log n\rceil\le2\log n$ 且 $4\le4\log n$，因此
\[
 L\le24n\log n\le\lfloor1000n\log n\rfloor.
\]
最后在末尾补零，使长度恰好为题目要求的长度。

\paragraph{单射性}
已知 $n$，可以按固定长度逐行解码，读出 $d_v$，再读出前 $d_v$ 个邻居。
这样能唯一恢复每个顶点的出邻接表，因而唯一恢复所有有向边。
所以两个图编码相同必然是同一个图，$E$ 为单射。

\section{习题 2.5(1)：把排列编码成有界度图}

\paragraph{题意}
令 $S_n$ 为所有双射 $\pi:[n]\to[n]$ 的集合。构造单射 $S_n\to G_{2n}$。

\paragraph{构造与证明}
把 $[2n]$ 分成左右两部分
\[
 L=\{0,\ldots,n-1\},\qquad R=\{n,\ldots,2n-1\}.
\]
对排列 $\pi$，定义有向图 $H_\pi$ 的边集为
\[
 E(H_\pi)=\bigl\{(i,n+\pi(i)):i\in[n]\bigr\}.
\]
每个左侧顶点恰有一条出边；因为 $\pi$ 是双射，每个右侧顶点恰有一条入边。
各顶点总度数都是 $1$，且两侧不相交，所以没有自环。因此 $H_\pi\in G_{2n}$。

若 $\pi\ne\sigma$，存在 $i$ 使得 $\pi(i)\ne\sigma(i)$，
于是从顶点 $i$ 出发的边在两图中不同，故 $H_\pi\ne H_\sigma$。
这就证明了单射性。

\paragraph{容易忽略的条件}
这里的图是顶点带有固定编号的图，并没有对图同构取商。
如果抹去编号，这些完美匹配会彼此同构，便不能区分排列。

\section{习题 2.14：有限整数序列的自然数编码}

\paragraph{题意}
构造单射 $\mathbb Z^*\to\mathbb N$，其中 $\mathbb Z^*$ 包含任意长度的有限整数序列及空序列。

\paragraph{第一步：编码单个整数}
定义双射 $h:\mathbb Z\to\mathbb N$（这里使用非负整数）为
\[
 h(z)=
 \begin{cases}
  2z,&z\ge0,\\
  -2z-1,&z<0.
 \end{cases}
\]
其逆映射为
\[
 h^{-1}(a)=
 \begin{cases}
  a/2,&a\text{ 为偶数},\\
  -(a+1)/2,&a\text{ 为奇数}.
 \end{cases}
\]

\paragraph{第二步：用素数指数保存位置与内容}
令 $p_0=2,p_1=3,p_2=5,\ldots$ 为递增的素数序列，定义
\[
 E(z_0,\ldots,z_{k-1})=\prod_{i=0}^{k-1}p_i^{h(z_i)+1},
 \qquad E(\varepsilon)=1.
\]
例如
\[
 E(-1,0,2)=2^2\cdot3^1\cdot5^5.
\]

\paragraph{单射性}
根据算术基本定理，编码整数的素因数分解唯一。
因为每个使用到的指数都至少是 $1$，非空序列的长度就是出现的连续前若干个素数的个数。
各指数减去 $1$ 后得到 $h(z_i)$，再用 $h^{-1}$ 得到 $z_i$。
数 $1$ 没有素因子，仅对应空序列。因此编码可以唯一解码。

这里的 ``$+1$'' 很重要：如果直接使用 $p_i^{h(z_i)}$，末尾的整数 $0$
会产生指数 $0$，从而无法区分一个序列与在其末尾添加 $0$ 后的序列。

\section{习题 3.9：LOOKUP 的通用性}

\paragraph{题意}
证明 $\{\LOOKUP_1,0,1\}$ 是通用门集合。教材的参数顺序为
\[
 \LOOKUP_1(a,b,c)=
 \begin{cases}
  a,&c=0,\\
  b,&c=1.
 \end{cases}
\]
也就是说，\textbf{第三个参数才是选择位}。

\paragraph{构造基本逻辑门}
直接检查可得
\begin{align*}
 \neg x&=\LOOKUP_1(1,0,x),\\
 x\land y&=\LOOKUP_1(0,y,x),\\
 x\lor y&=\LOOKUP_1(y,1,x).
\end{align*}
例如第二式中：$x=0$ 时输出 $0$；$x=1$ 时输出 $y$，正好是 $x\land y$。

由于 $\{\mathrm{AND},\mathrm{OR},\mathrm{NOT}\}$ 是通用的，
把任意布尔电路中的每个基本门替换成上述构造，就能得到仅使用
$\LOOKUP_1$ 和常数的等价电路。

也可以直接构造 NAND：先令 $t=\LOOKUP_1(1,0,y)=\neg y$，再令
\[
 \LOOKUP_1(1,t,x)=\neg(x\land y).
\]
所以每个 NAND 门至多用两个 LOOKUP 门替代，结论同样成立。

\section{习题 3.14：线性大小的多数电路}

\paragraph{题意}
证明存在绝对常数 $c$，使得每个 $n>1$ 的严格多数函数
\[
 \MAJ_n(x)=1\quad\Longleftrightarrow\quad
 \sum_{i=0}^{n-1}x_i>n/2
\]
都有不超过 $cn$ 个 NAND 门的电路。

\subsection*{1. 加法器的代价}
两个 $r$ 位二进制数可以用 $O(r)$ 个 NAND 门相加。
为说明这一点，一个全加器在一位上计算
\[
 p=a\oplus b,\qquad d=p\oplus c_{\mathrm{in}},\qquad
 c_{\mathrm{out}}=(a\land b)\lor(p\land c_{\mathrm{in}}).
\]
XOR 可用 $4$ 个 NAND 门实现，AND 用 $2$ 个，OR 用 $3$ 个，
因而上述一位全加器至多用 $15$ 个 NAND 门。
把 $r$ 个全加器按进位串联，就得到门数至多 $15r$ 的加法器，
输出允许有 $r+1$ 位。

若模型没有预置常数，可以从任意输入 $x_0$ 构造
\[
 u=\NAND(x_0,x_0),\quad
 \mathbf 1=\NAND(x_0,u),\quad
 \mathbf 0=\NAND(\mathbf 1,\mathbf 1),
\]
只需 $3$ 个门，随后共享这些常数线。

\subsection*{2. 平衡树中逐层相加}
令 $N=2^{\lceil\log n\rceil}$，在原输入后补 $N-n$ 个零。
于是 $n\le N<2n$，且补零不改变输入中 $1$ 的个数。

在第 $j$ 层，把两个大小为 $2^{j-1}$ 的块的计数相加，
得到大小为 $2^j$ 的块的计数，其中 $1\le j\le\log N$。
每个子块的计数至多为 $2^{j-1}$，用 $j$ 位足够，
所以这一层的每个加法器只需要 $O(j)$ 个门。
这一层共有 $N/2^j$ 个加法器，因此总门数至多
\[
 15\sum_{j=1}^{\log N}\frac{N}{2^j}j
 \le15N\sum_{j=1}^{\infty}\frac{j}{2^j}
 =30N<60n.
\]
最后得到二进制计数
\[
 S=\sum_{i=0}^{n-1}x_i,
\]
它只需要 $\lceil\log(n+1)\rceil$ 位，或保留树中多出的高位零。

这也对应递推式 $T(N)=2T(N/2)+O(\log N)$。
关键在于低层加法器处理的数很短，而不是让每次相加都使用 $\Theta(\log n)$ 位。

\subsection*{3. 与正确的阈值比较}
设置常数
\[
 t=\lfloor n/2\rfloor+1.
\]
对 $S$ 和 $t$ 从最高位向最低位扫描，维护两个状态：
$e$ 表示此前所有位都相等，$g$ 表示此前已经确定 $S>t$。
初始 $e=1,g=0$；读到位 $S_i,t_i$ 时同时更新为
\[
 g'=g\lor(e\land S_i\land\neg t_i),\qquad
 e'=e\land\neg(S_i\oplus t_i).
\]
每位只用常数个 NAND 门，故总成本为 $O(\log n)$。
最终输出 $g\lor e$，恰好判断 $S\ge t$。

所以总门数为 $O(n)+O(\log n)+O(1)=O(n)$，存在题目所需的常数 $c$。
\textbf{注意：}补零后比较的仍是原来的阈值 $n/2$，而不是 $N/2$；
当 $n$ 为偶数且恰有一半输入为 $1$ 时，输出应为 $0$。

\section{习题 5.1：哪一个求值器大小断言是假的？}

\paragraph{题意与答案}
输入为一个 NAND-CIRC 程序 $P$ 的编码，输出 $P$ 在全 $1$ 输入上的结果。
下列三项中，\textbf{第三项 (c) 为假}：
\begin{enumerate}[label=(\alph*)]
 \item $P$ 有 $s$ 行，使用 list-of-tuples 编码，求值器大小为 $O(s^3)$：真。
 \item $P$ 有 $s$ 个字符，使用 $7s$ 位 ASCII 编码，求值器大小为 $O(s^3)$：真。
 \item $P$ 有 $s$ 行，使用 list-of-tuples 编码，求值器大小为 $O(\sqrt s)$：假。
\end{enumerate}
这些断言理解为对合法程序编码正确；非法编码上的行为可以任意指定。
求值器可以依赖长度参数，但必须能处理该长度的所有合法程序。

\subsection*{(a) 元组编码的多项式大小求值器}
一个 $s$ 行程序只涉及 $O(s)$ 个实际出现的变量，
按教材规则重新编号后，每个地址只需 $O(\log s)$ 位。
模拟每一行时：
\begin{enumerate}
 \item 将两个源地址分别与 $O(s)$ 个存储位置的地址比较；
 \item 用比较结果选出两个当前值，执行一次 NAND；
 \item 对目标地址对应的位置写入结果，其余位置保留旧值。
\end{enumerate}
地址相等判断需 $O(\log s)$ 个门，单行模拟成本为 $O(s\log s)$。
所有输入变量的初值固定为 $1$，并不需要作为额外输入传给求值器。
将 $s$ 步模拟展开成一个电路，大小为
\[
 O(s^2\log s)\subseteq O(s^3).
\]
这里即使源程序重复写入同一变量，逐步更新存储的模拟仍然适用。

\subsection*{(b) 字符编码也有立方上界}
不能仅凭 ``解析是多项式时间'' 就直接得出立方上界。
下面给出一个较粗但足够的 $O(s^3)$ 门数构造。

把源文件中每一个字符位置都看作一个潜在的语句起点。
固定语法的分隔符、换行和变量名边界可通过字符测试及区间前缀电路确定。
对每个位置，提取该位置若为语句起点时的三个变量名：目标、第一源、第二源；
若不是起点，则将该槽标记为无效。每个名字统一补齐到 $O(s)$ 位，包含终止标记。
字符只有 $7$ 位，不能把字符数误当成比特数上的额外渐近因子。

提取这些字段时，对每个潜在语句、每个名字中的字符位置，
用至多 $s$ 个候选起点作选择，成本至多为 $O(s^3)$。
一共只有 $O(s)$ 个字段；对每一对字段比较名字是否相同，
每次用 $O(s)$ 个门，总成本仍为 $O(s^3)$。
不必把很长的下标展开成一个巨大数组；比较变量标识符即可。
如果语法允许数字下标有前导零，可以先去除前导零再比较，仍不改变该上界。

随后按字符位置从小到大处理有效语句。对一个源变量，
从之前语句中选出\emph{最后一次}写入同名变量的语句，取其已计算的值；
如果该名字是输入变量，则值固定为 $1$。
对固定语句，利用后缀 OR 得到 ``最后一次匹配'' 标志只需 $O(s)$ 个门，
选值和执行 NAND 也只需 $O(s)$ 个门。
展开所有语句需要 $O(s^2)$ 个门，最后选择输出变量的最终值即可。

因此字段提取、名字比较、求值三部分合计为 $O(s^3)$，(b) 成立。
上述构造只追求题目要求的上界，并不声称它最优。

\subsection*{(c) 一个严格的 $\Omega(s)$ 下界}
固定 $t\ge2$，构造由参数 $b=(b_1,\ldots,b_t)\in\bits^t$ 决定的程序 $P_b$。
它只有一个输入 $V_0$，求值时设为 $1$。先执行
\[
 V_1=\NAND(V_0,V_0)=0.
\]
再对每个 $i$ 新建一个变量 $B_i$，执行
\[
 B_i=\NAND(V_{b_i},V_{b_i}).
\]
于是 $b_i=0$ 时选取 $V_0=1$，得到 $B_i=0$；
$b_i=1$ 时选取 $V_1=0$，得到 $B_i=1$。
因此 $B_i=b_i$。

最后用两个 NAND 门实现一次 AND，依次求出
\[
 B_1\land B_2\land\cdots\land B_t.
\]
程序行数恰为
\[
 s=1+t+2(t-1)=3t-1,
\]
而它在全 $1$ 输入上的输出正是 $\bigwedge_{i=1}^t b_i$。

在元组编码中，改变 $b_i$ 只改变第 $i$ 个选值语句的两个源地址，
各 $i$ 对应的编码位置组彼此不交，所有变量编号、程序长度和其他元组均不变。
比较 $b=(1,\ldots,1)$ 与仅将 $b_i$ 改成 $0$ 的参数，程序输出由 $1$ 变为 $0$。
因此求值器必须使用第 $i$ 个位置组中的至少一个输入位，否则无法区分这两个合法编码。

求值器于是必须依赖至少 $t$ 个不同的输入位。
一个具有 $q$ 个二输入 NAND 门的非平凡单输出电路，最多直接使用 $2q$ 个不同输入位，故
\[
 2q\ge t,\qquad q\ge t/2=\Omega(s).
\]
这不可能同时是 $O(\sqrt s)$，所以 (c) 为假。
这里不是笼统地声称 ``必须读完整个程序''，而是展示了必须区分的合法程序族。

\section{习题 5.2：比较两个输入字符串}

\paragraph{题意}
构造 $O(k)$ 行的 NAND-CIRC 程序，计算
\[
 \mathrm{EQUALS}_k(x,y)=1\quad\Longleftrightarrow\quad x=y,
 \qquad x,y\in\bits^k.
\]

\paragraph{逐位判断相等}
对 $i=0,\ldots,k-1$，执行
\begin{align*}
 a_i&=\NAND(x_i,y_i),\\
 b_i&=\NAND(x_i,a_i),\\
 c_i&=\NAND(y_i,a_i),\\
 d_i&=\NAND(b_i,c_i)=x_i\oplus y_i,\\
 e_i&=\NAND(d_i,d_i)=\neg(x_i\oplus y_i).
\end{align*}
每一位恰用 $5$ 个门，且 $e_i=1$ 当且仅当对应的两个输入位相等。

\paragraph{合并所有位}
令 $r_0=e_0$，对 $i=1,\ldots,k-1$ 执行
\[
 u_i=\NAND(r_{i-1},e_i),\qquad r_i=\NAND(u_i,u_i).
\]
于是 $r_i=r_{i-1}\land e_i$，最终输出为
\[
 r_{k-1}=\bigwedge_{i=0}^{k-1}\neg(x_i\oplus y_i).
\]
当 $k\ge1$ 时，总门数为
\[
 5k+2(k-1)=7k-2=O(k).
\]
这两个循环是在\emph{构造程序时}展开的；得到的 NAND-CIRC 程序本身没有循环。
最后一个门的目标变量直接命名为输出即可。

\paragraph{空字符串情形}
若题目的自然数包含 $0$，则 $k=0$ 时两个输入都是空字符串，函数恒为 $1$。
在允许常数输出的约定下单独处理即可；通常的渐近 $O(k)$ 结论针对 $k\ge1$。

\section{习题 5.3：与一个固定字符串比较}

\paragraph{题意}
给定固定的 $x'\in\bits^k$，构造 $O(k)$ 行的 NAND-CIRC 程序计算
\[
 \mathrm{EQUALS}_{x'}(x)=1\quad\Longleftrightarrow\quad x=x'.
\]
这里 $x'$ 是构造程序时已知的常量，不是运行时的另一组输入。

\paragraph{构造}
对每一位定义文字
\[
 \ell_i=
 \begin{cases}
  x_i,&x'_i=1,\\
  \neg x_i,&x'_i=0.
 \end{cases}
\]
显然 $\ell_i=1$ 当且仅当 $x_i=x'_i$，所以
\[
 \mathrm{EQUALS}_{x'}(x)=\bigwedge_{i=0}^{k-1}\ell_i.
\]
常量位为 $1$ 时直接使用输入线；常量位为 $0$ 时使用一行
$\ell_i=\NAND(x_i,x_i)$。
再按习题 5.2 的方法，用 $2(k-1)$ 个门求所有文字的 AND。

若 $x'$ 有 $z$ 个零，当 $k\ge2$ 时门数为
\[
 z+2(k-1)\le3k-2.
\]
若 $k=1$ 且目标输出就是输入线，而语法要求输出必须由赋值产生，
可用两次取反复制输入；$k=1$ 且需要取反则只需一个门。
因此统一有 $O(k)$ 上界。$k=0$ 的约定与上一题相同。

\section{习题 5.4：多输出电路下界与原题的量词问题}

\paragraph{原题}
原文要求证明存在常数 $\delta>0$，使得对每个充分大的 $n$ 和\emph{每个} $m$，
都存在函数 $f:\bits^n\to\bits^m$，任何计算 $f$ 的 NAND 电路都至少需要
\[
 \delta\frac{m2^n}{n}
\]
个门。

\subsection*{1. 为什么不限 $m$ 时不成立}
固定 $n$。所有单输出函数 $g:\bits^n\to\bits$ 只有
\[
 Q_n=2^{2^n}
\]
个。利用析取范式，每个这样的函数都可以由至多 $B_n=O(n2^n)$ 个 NAND 门计算。
把这些电路全部放在一起，共享同一组 $n$ 个输入，便得到一个计算\emph{所有}
单输出函数的电路库，其总门数不超过某个整数
\[
 A_n\le Q_nB_n.
\]
这里只需知道 $A_n$ 是一个依赖于 $n$ 的有限整数，并不需要它很小。

任意 $m$ 输出函数都可写成
\[
 f(x)=(f_1(x),\ldots,f_m(x)),
\]
每个 $f_j$ 都已经由电路库中的某一条线算出。
如果模型允许重复使用输出线，只需把这些线接到输出端；
即使要求每个输出对应不同的门，也可对每个输出附加两个 NAND 门：
\[
 a_j=\NAND(f_j,f_j),\qquad Y_j=\NAND(a_j,a_j)=f_j.
\]
因此\textbf{每一个} $m$ 输出函数都有大小至多 $A_n+2m$ 的电路。

现在假设原题中的 $\delta>0$ 存在。选取足够大的 $n$，使
\[
 \delta\frac{2^n}{n}>3,
\]
再选取整数 $m\ge A_n$。此时所有 $m$ 输出函数都满足
\[
 \text{最小电路大小}\le A_n+2m\le3m
 <\delta\frac{m2^n}{n}.
\]
这与原题声称存在某个函数达到该下界矛盾。

因此，\textbf{原文 ``对每个 $m$'' 的版本是假的}。
直观原因是：当输出数量极大时，可以预先算出所有可能的单输出函数，
后面的输出只需要重复使用或复制这些结果。

\subsection*{2. 一个正确版本及完整计数证明}
一个常用且足够明确的修正版是：存在绝对常数 $\delta>0$，
使得对充分大的 $n$ 及 $1\le m\le2^n$，存在 $f:\bits^n\to\bits^m$，其电路大小至少为
$\delta m2^n/n$。

\paragraph{函数总数}
输入共有 $2^n$ 种，每个输入的输出有 $2^m$ 种选择，故函数总数为
\[
 (2^m)^{2^n}=2^{m2^n}.
\]

\paragraph{小电路的数量}
考虑 $s\ge\max\{n,m,2\}$。将任意至多 $s$ 个门的电路补上无用门，
并按拓扑序排列为恰好 $s$ 个门。
每个 NAND 门的两个源各有至多 $n+s$ 种选择，
每个输出端再从至多 $n+s$ 条线中选一个。
允许不合法的选项只会扩大计数，故不同函数的数量至多
\[
 (n+s)^{2s+m}\le(2s)^{3s}
 \le 2^{10s\log s}.
\]
这里不同编码可能计算相同函数，所以这是上界，不必要求编码唯一。

\paragraph{选择参数}
取 $\delta=1/40$，令
\[
 s=\left\lfloor\delta\frac{m2^n}{n}\right\rfloor.
\]
当 $n$ 充分大时，对所有 $1\le m\le2^n$，都有 $s\ge\max\{n,m,2\}$。
另一方面，$s\le m2^n\le2^{2n}$，所以 $\log s\le2n$。于是
\[
 10s\log s
 \le20\delta m2^n
 =\frac12m2^n
 <m2^n.
\]
因此至多 $s$ 个门的电路所能计算的函数总数，严格少于所有函数总数。
必有某个函数不能由至多 $s$ 个门的电路计算。
它的最小门数是整数，故至少为
\[
 s+1>\delta\frac{m2^n}{n},
\]
得到所需下界。

\paragraph{证明中究竟哪里需要限制 $m$}
关键是 $\log s=O(n)$。
如果任意放大 $m$，一般只能得到 $\log s=O(n+\log m)$，
不能把这个因子无条件替换为 $O(n)$。
同样的计数方法可推广到任意固定 $C>0$ 的范围 $m\le2^{Cn}$，
但下界中的常数需要允许依赖 $C$。

\section{习题 5.9：由元组恢复程序和函数}

\paragraph{题目给出的编码}
输入数 $n=3$，输出数 $m=1$，元组依次为
\[
\begin{gathered}
 (3,2,2),\ (4,1,1),\ (5,3,4),\ (6,2,1),\ (7,6,6),\\
 (8,0,0),\ (9,7,8),\ (10,5,0),\ (11,9,10).
\end{gathered}
\]
每个三元组 $(i,j,k)$ 表示 $V_i=\NAND(V_j,V_k)$。
按本书第 5.3 节的编号约定，最前面的 $n$ 个变量是输入，
\textbf{最后面的 $m$ 个变量是输出}。因此本题输出为 $V_{11}$，不是 $V_3$。

\paragraph{逐行求值}
记 $a=x_0=V_0$、$b=x_1=V_1$、$c=x_2=V_2$。则
\[
\begin{array}{ccl}
\toprule
\text{变量}&\text{程序语句}&\text{化简后的值}\\
\midrule
V_3&\NAND(V_2,V_2)&\neg c\\
V_4&\NAND(V_1,V_1)&\neg b\\
V_5&\NAND(V_3,V_4)&b\lor c\\
V_6&\NAND(V_2,V_1)&\neg(b\land c)\\
V_7&\NAND(V_6,V_6)&b\land c\\
V_8&\NAND(V_0,V_0)&\neg a\\
V_9&\NAND(V_7,V_8)&\neg(b\land c\land\neg a)\\
V_{10}&\NAND(V_5,V_0)&\neg((b\lor c)\land a)\\
V_{11}&\NAND(V_9,V_{10})&(b\land c\land\neg a)\lor(a\land(b\lor c))\\
\bottomrule
\end{array}
\]

\paragraph{化简输出}
由德摩根律和分配律，输出为
\[
 f(a,b,c)=(\neg a\land b\land c)\lor(a\land b)\lor(a\land c).
\]
再利用 $X\lor(\neg a\land Y)=(X\lor\neg a)\land(X\lor Y)$，
或直接对 $a$ 分类，可得
\[
 \boxed{f(a,b,c)=(a\land b)\lor(a\land c)\lor(b\land c)=\MAJ_3(a,b,c).}
\]
更直观地说：$a=0$ 时必须 $b=c=1$ 才输出 $1$；
$a=1$ 时只需 $b,c$ 至少一个为 $1$。两种情况合在一起，就是三位中至少两位为 $1$。

\paragraph{真值表}
按输入 $000,001,010,011,100,101,110,111$ 的字典序：
\[
\begin{array}{ccc|c}
 a&b&c&f(a,b,c)\\
\hline
 0&0&0&0\\
 0&0&1&0\\
 0&1&0&0\\
 0&1&1&1\\
 1&0&0&0\\
 1&0&1&1\\
 1&1&0&1\\
 1&1&1&1
\end{array}
\]
因此题目所问的 $8$ 元组为
\[
 \boxed{(0,0,0,1,0,1,1,1).}
\]

\clearpage
\phantomsection
\addcontentsline{toc}{part}{English Solutions}
\part*{English Solutions}

\noindent
Throughout this part, all logarithms are base $2$, and
$[n]=\{0,\ldots,n-1\}$. Circuit size means the number of NAND gates;
wires and fan-out do not contribute to the gate count. Statements about
evaluators require correctness on valid program encodings.

\section{Exercise 2.4: Encoding Bounded-Degree Graphs}

\paragraph{Claim.}
Let $G_n$ be the set of directed graphs on the labeled vertex set $[n]$
with no self-loops and degree at most $10$ at every vertex.
For all sufficiently large $n$, there is an injection
\[
 E:G_n\longrightarrow\bits^{\lfloor1000n\log n\rfloor}.
\]

\paragraph{Construction.}
Set $\ell=\lceil\log n\rceil$. For each vertex $v$, in increasing order,
record its outgoing adjacency list as follows:
\begin{enumerate}
 \item Store its out-degree $d_v\in\{0,\ldots,10\}$ using $4$ bits.
 \item List its outgoing neighbors in increasing order, using $\ell$ bits
       per vertex label and reserving exactly $10$ slots.
 \item Fill any unused slots with zeros.
\end{enumerate}
The out-degree is at most $10$ whether the degree bound is interpreted
as a bound on total degree or on out-degree. A real neighbor labeled $0$
cannot be confused with padding because $d_v$ specifies the number of
occupied slots.

The resulting string has length
\[
 L=n\bigl(4+10\lceil\log n\rceil\bigr).
\]
For $n\ge2$, we have $\lceil\log n\rceil\le2\log n$ and
$4\le4\log n$, so
\[
 L\le24n\log n\le\lfloor1000n\log n\rfloor.
\]
Append zeros to obtain exactly the required length.

\paragraph{Injectivity.}
Since $n$ is known, the string can be split into fixed-length records.
The decoder reads $d_v$ and the first $d_v$ neighbor labels in each record.
This uniquely recovers every outgoing adjacency list and hence every edge.
Consequently, $E(G)=E(H)$ implies $G=H$, proving that $E$ is injective.

\section{Exercise 2.5(1): Encoding Permutations as Graphs}

\paragraph{Claim.}
If $S_n$ is the set of bijections $\pi:[n]\to[n]$, then there is an
injection $S_n\to G_{2n}$.

\paragraph{Proof.}
Partition the labeled vertex set $[2n]$ into
\[
 L=\{0,\ldots,n-1\},\qquad R=\{n,\ldots,2n-1\}.
\]
For $\pi\in S_n$, define $H_\pi$ by
\[
 E(H_\pi)=\{(i,n+\pi(i)):i\in[n]\}.
\]
Every vertex in $L$ has exactly one outgoing edge. Since $\pi$ is a
bijection, every vertex in $R$ has exactly one incoming edge.
Thus every vertex has total degree $1$, and there are no self-loops
because the two parts are disjoint. In particular, $H_\pi\in G_{2n}$.

If $\pi\ne\sigma$, choose $i$ with $\pi(i)\ne\sigma(i)$.
The unique edge leaving $i$ is different in $H_\pi$ and $H_\sigma$,
so these graphs are distinct. Therefore $\pi\mapsto H_\pi$ is injective.

The fixed vertex labels are essential: the exercise concerns labeled
graphs, not isomorphism classes. Without labels, all these matchings
would be isomorphic.

\section{Exercise 2.14: Encoding Finite Integer Sequences}

\paragraph{Claim.}
There is an injection $\mathbb Z^*\to\mathbb N$, including an encoding
of the empty sequence.

\paragraph{Construction.}
First encode each integer by the nonnegative integer
\[
 h(z)=
 \begin{cases}
  2z,&z\ge0,\\
  -2z-1,&z<0.
 \end{cases}
\]
This is a bijection: an even value $a$ decodes to $a/2$, and an odd
value $a$ decodes to $-(a+1)/2$.

Let $p_0=2,p_1=3,p_2=5,\ldots$ be the primes in increasing order.
Define
\[
 E(z_0,\ldots,z_{k-1})=\prod_{i=0}^{k-1}p_i^{h(z_i)+1},
 \qquad E(\varepsilon)=1.
\]
For example, $E(-1,0,2)=2^2\cdot3\cdot5^5$.

\paragraph{Injectivity.}
The fundamental theorem of arithmetic gives a unique prime
factorization of the encoded integer. All exponents in the product
are positive, so the prime factors are exactly the first $k$ primes.
We can therefore recover the length $k$, subtract $1$ from each exponent,
and apply $h^{-1}$ to recover the entries in order.
Only the empty sequence maps to $1$.
Thus every encoded sequence is uniquely recoverable, proving injectivity.

Adding $1$ to each exponent is necessary for this construction:
without it, appending a zero to a sequence would introduce an exponent
of zero and would not change the encoded number.

\section{Exercise 3.9: Universality of LOOKUP}

\paragraph{Claim.}
The gate set $\{\LOOKUP_1,0,1\}$ is universal, where
\[
 \LOOKUP_1(a,b,c)=
 \begin{cases}
  a,&c=0,\\
  b,&c=1.
 \end{cases}
\]
The selector is the \emph{third} argument.

\paragraph{Proof.}
The usual Boolean gates have the following implementations:
\begin{align*}
 \neg x&=\LOOKUP_1(1,0,x),\\
 x\land y&=\LOOKUP_1(0,y,x),\\
 x\lor y&=\LOOKUP_1(y,1,x).
\end{align*}
Each identity follows by considering $x=0$ and $x=1$.
Since AND, OR, and NOT form a universal basis, replacing their
occurrences by these constructions proves the claim.

Equivalently, set $t=\LOOKUP_1(1,0,y)=\neg y$. Then
\[
 \LOOKUP_1(1,t,x)=\neg(x\land y).
\]
Thus two LOOKUP gates implement one NAND gate, and universality
also follows directly from the universality of NAND.

\section{Exercise 3.14: A Linear-Size Majority Circuit}

\paragraph{Claim.}
There is an absolute constant $c$ such that, for every $n>1$,
the strict majority function
\[
 \MAJ_n(x)=1\quad\Longleftrightarrow\quad
 \sum_{i=0}^{n-1}x_i>n/2
\]
can be computed using at most $cn$ NAND gates.

\subsection*{Step 1: Linear-size binary adders}
A full adder computes
\[
 p=a\oplus b,\qquad d=p\oplus c_{\mathrm{in}},\qquad
 c_{\mathrm{out}}=(a\land b)\lor(p\land c_{\mathrm{in}}).
\]
XOR, AND, and OR use at most $4$, $2$, and $3$ NAND gates respectively.
Hence a full adder uses at most $15$ NAND gates. A ripple-carry adder
for two $r$-bit numbers uses at most $15r$ gates and produces $r+1$
output bits, including the final carry.

If constant wires are not available initially, construct them once:
\[
 u=\NAND(x_0,x_0),\quad
 \mathbf 1=\NAND(x_0,u),\quad
 \mathbf 0=\NAND(\mathbf 1,\mathbf 1).
\]
This costs three gates, after which the constant wires may be shared.

\subsection*{Step 2: Add the input bits in a balanced tree}
Let $N=2^{\lceil\log n\rceil}$ and append $N-n$ zero bits.
Then $n\le N<2n$, and the number of ones is unchanged.

At level $j$, for $1\le j\le\log N$, combine two counts for blocks of
$2^{j-1}$ input bits. Each child count is at most $2^{j-1}$ and therefore
fits in $j$ bits. Thus each addition at this level uses at most $15j$
gates. There are $N/2^j$ additions at level $j$, so all additions together
use at most
\[
 15\sum_{j=1}^{\log N}\frac{N}{2^j}j
 \le15N\sum_{j=1}^{\infty}\frac{j}{2^j}
 =30N<60n
\]
gates. The result is the binary representation of
$S=\sum_{i=0}^{n-1}x_i$, using $O(\log n)$ bits.

The corresponding recurrence is $T(N)=2T(N/2)+O(\log N)$.
The linear bound arises because additions near the leaves involve
short numbers; charging $\Theta(\log n)$ for every addition would
give an unnecessarily large upper bound.

\subsection*{Step 3: Compare against the threshold}
Set $t=\lfloor n/2\rfloor+1$. Compare $S$ with $t$ from most significant
bit to least significant bit. Maintain bits $e$ and $g$, where $e$
means the prefixes seen so far are equal and $g$ means the prefix of
$S$ is larger. Initially $e=1$ and $g=0$. At position $i$, update
simultaneously by
\[
 g'=g\lor(e\land S_i\land\neg t_i),\qquad
 e'=e\land\neg(S_i\oplus t_i).
\]
Each position requires only a constant number of NAND gates.
At the end, $g\lor e$ is $1$ exactly when $S\ge t$, so the comparator
costs $O(\log n)$ gates.

The total size is $O(n)+O(\log n)+O(1)=O(n)$, as required.
The threshold is determined by the \emph{original} length $n$, not the
padded length $N$. In particular, a tie when $n$ is even must produce $0$.

\section{Exercise 5.1: Which Evaluator Claim Is False?}

\paragraph{Answer.}
\textbf{Statement (c) is false.} Statements (a) and (b) are true.
The evaluator receives the description of a NAND-CIRC program and
must compute its output when all of the program's inputs are $1$.

\subsection*{(a) Tuple encoding: the cubic bound is valid}
A program of $s$ lines mentions at most $O(s)$ distinct variables.
With the textbook's renumbering convention, each address has
$O(\log s)$ bits.

For each line, compare its two source addresses against all $O(s)$
memory addresses, select the corresponding current values, and apply
NAND. Then compare the destination address against the memory
addresses and update the selected location, retaining every other
value. Each address comparison uses $O(\log s)$ gates, so one step
costs $O(s\log s)$ gates.

Set all input locations initially to $1$. Unrolling all $s$ steps gives
an evaluator of size
\[
 O(s^2\log s)\subseteq O(s^3).
\]
This construction also handles repeated assignments to a variable.

\subsection*{(b) ASCII encoding: the cubic bound is also valid}
Here $s$ counts characters, and the input encoding contains $7s$ bits.
Simply saying that parsing is polynomial-time would not establish
the particular cubic bound, so we describe a sufficiently conservative
circuit construction.

Treat every character position as a possible statement start.
Character tests and prefix computations identify separators, line
boundaries, and variable-name boundaries in the fixed syntax.
For each possible start, extract the destination and the two source
identifiers, padding each to length $O(s)$ with an end marker.
Mark a slot invalid if it does not begin a statement.

There are $O(s)$ fields. For each character in each field, selecting
among at most $s$ possible starting positions costs $O(s)$ gates;
thus extracting all fields costs at most $O(s^3)$ gates.
Compare every pair of identifiers using $O(s)$ gates per comparison.
There are $O(s^2)$ pairs, so this also costs $O(s^3)$ gates.
The fixed seven-bit character width contributes only a constant factor.
Large numerical subscripts need not be expanded into a large memory
array: we compare identifiers instead. If leading zeros are permitted
in subscripts, normalize them first, within the same bound.

Now evaluate valid statement slots in their original order. For each
source identifier, use its comparisons with earlier destinations to
select the value from the \emph{last} matching assignment; input
identifiers instead have value $1$. For a fixed source, the last-match
flags are obtained by a suffix-OR scan in $O(s)$ gates. Selecting the
value and computing the new NAND result also costs $O(s)$ gates per
statement. Unrolling this stage costs $O(s^2)$ gates. Finally select
the final values of the output variables.

The total cost is therefore $O(s^3)$, proving (b).

\subsection*{(c) A linear lower bound rules out the square-root bound}
For each $t\ge2$ and $b=(b_1,\ldots,b_t)\in\bits^t$, construct a
single-output program $P_b$ with one input $V_0$. On the required
all-ones input, $V_0=1$. Begin with
\[
 V_1=\NAND(V_0,V_0)=0.
\]
For each $i$, create a fresh variable $B_i$ by the instruction
\[
 B_i=\NAND(V_{b_i},V_{b_i}).
\]
If $b_i=0$, this negates $V_0=1$ and produces $0$; if $b_i=1$, it
negates $V_1=0$ and produces $1$. Hence $B_i=b_i$.

Finally compute $B_1\land\cdots\land B_t$, using two NAND gates
for each of the $t-1$ binary AND operations. The number of lines is
\[
 s=1+t+2(t-1)=3t-1,
\]
and the program's output on the all-ones input is $\bigwedge_i b_i$.

Changing $b_i$ changes only the two source-address fields in its own
instruction. These groups of encoding positions are disjoint for
different $i$; the program length, variable numbering, and all other
tuples remain unchanged.

Compare the encoding for $b=(1,\ldots,1)$ with the encoding obtained
by flipping only $b_i$ to zero. Their required outputs differ.
Consequently, a correct evaluator must use at least one input bit
from the $i$th position group, for every $i$. It must therefore depend
on at least $t$ distinct input bits.

A nontrivial single-output circuit with $q$ two-input NAND gates has
only $2q$ gate-input ports, and hence can use at most $2q$ distinct
external input bits. It follows that
\[
 q\ge t/2=\Omega(s).
\]
This contradicts an $O(\sqrt{s})$ bound along the unbounded sequence
$s=3t-1$. Thus (c) is false. The argument uses only valid program
encodings, not an assumption that the evaluator must read every bit.

\section{Exercise 5.2: Equality of Two Strings}

\paragraph{Claim.}
For $k\ge1$, the function
\[
 \mathrm{EQUALS}_k(x,y)=1\quad\Longleftrightarrow\quad x=y,
 \qquad x,y\in\bits^k,
\]
has an $O(k)$-line NAND-CIRC program.

\paragraph{Construction.}
For each position $i$, compute
\begin{align*}
 a_i&=\NAND(x_i,y_i),\\
 b_i&=\NAND(x_i,a_i),\\
 c_i&=\NAND(y_i,a_i),\\
 d_i&=\NAND(b_i,c_i)=x_i\oplus y_i,\\
 e_i&=\NAND(d_i,d_i)=\neg(x_i\oplus y_i).
\end{align*}
These five gates make $e_i=1$ exactly when $x_i=y_i$.
Set $r_0=e_0$ and, for $i=1,\ldots,k-1$, compute
\[
 u_i=\NAND(r_{i-1},e_i),\qquad r_i=\NAND(u_i,u_i).
\]
Then $r_i=r_{i-1}\land e_i$, and the final output is
\[
 r_{k-1}=\bigwedge_{i=0}^{k-1}\neg(x_i\oplus y_i).
\]
It equals $1$ if and only if the entire strings agree.
The number of gates is exactly
\[
 5k+2(k-1)=7k-2=O(k).
\]
The indexed instructions are expanded when constructing the program;
the resulting NAND-CIRC program is straight-line code, without loops.
The last gate can be assigned directly to the output variable.

If $k=0$ is included, the only pair consists of two empty strings,
and equality is the constant-$1$ function. Handle this separately
under the convention allowing constant outputs; the asymptotic bound
above concerns positive $k$.

\section{Exercise 5.3: Equality to a Fixed String}

\paragraph{Claim.}
For each fixed $x'\in\bits^k$, the function
\[
 \mathrm{EQUALS}_{x'}(x)=1\quad\Longleftrightarrow\quad x=x'
\]
has an $O(k)$-line NAND-CIRC program.

\paragraph{Proof.}
The string $x'$ is known when the program is constructed; it is not
another runtime input. Define the literals
\[
 \ell_i=
 \begin{cases}
  x_i,&x'_i=1,\\
  \neg x_i,&x'_i=0.
 \end{cases}
\]
Each literal is $1$ exactly when its corresponding input bit matches
the fixed bit. Therefore
\[
 \mathrm{EQUALS}_{x'}(x)=\bigwedge_{i=0}^{k-1}\ell_i.
\]
Use the input wire directly when $x'_i=1$, and use one NAND gate
$\NAND(x_i,x_i)$ when $x'_i=0$. Combine the literals with $k-1$
AND operations, each implemented using two NAND gates.

If $x'$ contains $z$ zeros and $k\ge2$, the gate count is
\[
 z+2(k-1)\le3k-2.
\]
For $k=1$, a negation costs one gate. If the output must be an
assigned variable and the desired function is the identity, copy
the input using two successive negations. Thus the bound is $O(k)$
in all positive-length cases. The empty-string case is handled as
in Exercise 5.2.

\section{Exercise 5.4: A Counting Bound and a Missing Restriction}

\paragraph{The statement as printed.}
The exercise asserts that there is a constant $\delta>0$ such that,
for every sufficiently large $n$ and \emph{every} $m$, some function
$f:\bits^n\to\bits^m$ requires at least $\delta m2^n/n$ NAND gates.
\textbf{With no restriction on $m$, this statement is false.}

\subsection*{1. Counterexample to the unrestricted statement}
Fix $n$. There are exactly $Q_n=2^{2^n}$ single-output Boolean
functions on $n$ input bits. Each has a NAND circuit of size at most
$B_n=O(n2^n)$, for example by converting its disjunctive normal form
to NAND gates.

Place circuits for all these functions together, sharing their
$n$ input wires. This gives a finite circuit library of size at most
some integer $A_n\le Q_nB_n$, with a wire for every possible
single-output function. The important fact is that $A_n$ depends
only on $n$, not on the eventual number of outputs.

For any $f:\bits^n\to\bits^m$, write
\[
 f(x)=(f_1(x),\ldots,f_m(x)).
\]
Each coordinate function $f_j$ is already available in the library.
Even if the circuit definition requires distinct gates for distinct
outputs, two extra gates per coordinate suffice:
\[
 a_j=\NAND(f_j,f_j),\qquad Y_j=\NAND(a_j,a_j)=f_j.
\]
Thus \emph{every} such $f$ has a circuit of size at most $A_n+2m$.

Suppose the asserted constant $\delta>0$ existed. Choose $n$ large
enough both for the assertion to apply and for $\delta2^n/n>3$.
Then choose an integer $m\ge A_n$. For every function of these
input and output lengths, the preceding construction gives
\[
 \text{minimum circuit size}\le A_n+2m\le3m
 <\delta\frac{m2^n}{n}.
\]
This contradicts the claimed existence of a function requiring that
many gates. Hence the original statement cannot hold for all $m$.

The issue is sharing: with sufficiently many outputs, one may
precompute every possible coordinate function and then reuse it.

\subsection*{2. A corrected statement}
A valid version is the following: there is an absolute constant
$\delta>0$ such that, for every sufficiently large $n$ and every
$1\le m\le2^n$, some $f:\bits^n\to\bits^m$ requires at least
$\delta m2^n/n$ NAND gates. This restriction is supplied here; it
is not present in the printed exercise.

\paragraph{Counting all functions.}
There are $2^n$ inputs, and each input may independently receive any
of $2^m$ output strings. Hence the number of functions is
\[
 (2^m)^{2^n}=2^{m2^n}.
\]

\paragraph{Counting small circuits.}
Let $s\ge\max\{n,m,2\}$. Pad any circuit of at most $s$ gates
with unused gates and list its gates in topological order.
Each gate has two sources, each with at most $n+s$ choices.
Each of the $m$ output positions has at most $n+s$ choices of wire.
Allowing invalid descriptions only increases this count, so the
number of functions computed by these circuits is at most
\[
 (n+s)^{2s+m}\le(2s)^{3s}\le2^{10s\log s}.
\]
Distinct descriptions may compute the same function, which is
harmless because we need an upper bound.

\paragraph{Choosing the size threshold.}
Take $\delta=1/40$ and let
\[
 s=\left\lfloor\delta\frac{m2^n}{n}\right\rfloor.
\]
For all sufficiently large $n$, uniformly over $1\le m\le2^n$,
we have $s\ge\max\{n,m,2\}$. Also
$s\le m2^n\le2^{2n}$, so $\log s\le2n$. Consequently,
\[
 10s\log s\le20\delta m2^n=\frac12m2^n<m2^n.
\]
There are strictly fewer functions computed by circuits of at most
$s$ gates than there are functions altogether. Some function
therefore requires at least
\[
 s+1>\delta\frac{m2^n}{n}
\]
gates, proving the corrected claim.

\paragraph{Where the restriction is used.}
The crucial estimate is $\log s=O(n)$. With unrestricted $m$, the
corresponding estimate involves $n+\log m$, and this cannot be
replaced uniformly by $O(n)$.
The same counting proof works for $m\le2^{Cn}$ for any fixed $C>0$,
with a lower-bound constant allowed to depend on $C$.

\section{Exercise 5.9: Decoding the NAND Program}

\paragraph{Encoding convention.}
The program has $n=3$ inputs and $m=1$ output, with instructions
\[
\begin{gathered}
 (3,2,2),\ (4,1,1),\ (5,3,4),\ (6,2,1),\ (7,6,6),\\
 (8,0,0),\ (9,7,8),\ (10,5,0),\ (11,9,10).
\end{gathered}
\]
A tuple $(i,j,k)$ means $V_i=\NAND(V_j,V_k)$.
Under the textbook's convention, the first $n$ variables are inputs
and the \emph{last} $m$ variables are outputs. Thus the output is
$V_{11}$, not $V_3$.

\paragraph{Evaluating the instructions.}
Write $a=V_0$, $b=V_1$, and $c=V_2$. Successive simplification gives
\begin{align*}
 V_3&=\neg c,& V_4&=\neg b,\\
 V_5&=b\lor c,& V_6&=\neg(b\land c),\\
 V_7&=b\land c,& V_8&=\neg a,\\
 V_9&=\neg(b\land c\land\neg a),&
 V_{10}&=\neg((b\lor c)\land a).
\end{align*}
By De Morgan's law, the final output is
\begin{align*}
 f(a,b,c)=V_{11}
 &=(\neg a\land b\land c)\lor(a\land(b\lor c))\\
 &=(\neg a\land b\land c)\lor(a\land b)\lor(a\land c)\\
 &=(a\land b)\lor(a\land c)\lor(b\land c).
\end{align*}
For the last equality, combine
$(a\land b)\lor(\neg a\land b\land c)
=b\land(a\lor(\neg a\land c))=b\land(a\lor c)$.

\paragraph{1. Truth table.}
\[
\begin{array}{ccc|c}
 a&b&c&P(abc)\\
\hline
 0&0&0&0\\
 0&0&1&0\\
 0&1&0&0\\
 0&1&1&1\\
 1&0&0&0\\
 1&0&1&1\\
 1&1&0&1\\
 1&1&1&1
\end{array}
\]
In the requested order $000,001,010,011,100,101,110,111$, the output
tuple is
\[
 \boxed{(0,0,0,1,0,1,1,1).}
\]

\paragraph{2. Description in words.}
The program computes the majority function on three bits. It outputs
$1$ exactly when at least two of its three input bits are $1$.
Indeed, if $a=0$, both $b$ and $c$ must be $1$; if $a=1$, at least
one of $b,c$ must be $1$.

\end{document}
